\section{Objet de l'annexe}
\label{annexe-existence-objet}

Cette annexe justifie les propriétés élémentaires utilisées dans le
Tome I :

\begin{itemize}
\item
  les bornes associées à la matrice relationnelle, à la trace et au
  préconditionneur \(M(\tau)\) ;
\item
  l'existence et l'unicité de la dynamique continue pour un graphe
  fixé ;
\item
  l'absence d'explosion en temps fini ;
\item
  le caractère bien défini des sauts ;
\item
  l'absence d'une accumulation infinie de sauts en temps fini.
\end{itemize}

Le dernier point peut être établi sans supposer que les trajectoires
restent uniformément bornées lorsque \(t\to\infty\). Il suffit de
disposer d'une borne sur chaque intervalle de temps fini. Cette borne
est obtenue directement ci-dessous.

\section{Normes et état continu}
\label{annexe-existence-normes}

Le graphe \(A\) est d'abord fixé. L'état continu est

\[
z=(H,V,\tau,P),
\]

avec

\[
H,V,\tau\in\mathbb R^{N\times d},
\qquad
P\in\mathbb R_+^N.
\]

Pour une famille \(Y=(Y_1,\ldots,Y_N)\), avec
\(Y_i\in\mathbb R^d\), on utilise la norme

\[
\|Y\|_{\infty,2}
=
\max_{1\leq i\leq N}\|Y_i\|_2.
\]

Toutes les normes sont équivalentes en dimension finie. Ce choix est
simplement adapté à une matrice stochastique par lignes.

L'entrée extérieure est supposée mesurable et essentiellement bornée :

\[
\|u_i(t)\|_2\leq u_{\max}
\]

pour presque tout \(t\).

\section{Matrice relationnelle normalisée}
\label{annexe-existence-matrice-relationnelle}

On rappelle

\[
\widetilde A
=
\operatorname{Deg}(A+I)^{-1}(A+I).
\]

L'ajout de \(I\) garantit que chaque somme de ligne de \(A+I\) est
strictement positive. La matrice \(\operatorname{Deg}(A+I)\) est donc
inversible.

Pour tout \(i\),

\[
\widetilde A_{ij}\geq0,
\qquad
\sum_{j=1}^N\widetilde A_{ij}=1.
\]

Ainsi,

\[
(\widetilde A H)_i
=
\sum_{j=1}^N\widetilde A_{ij}H_j
\]

est une combinaison convexe des vecteurs \(H_j\). Par l'inégalité
triangulaire,

\[
\begin{aligned}
\|(\widetilde A H)_i\|_2
&\leq
\sum_j\widetilde A_{ij}\|H_j\|_2
\\
&\leq
\|H\|_{\infty,2}
\sum_j\widetilde A_{ij}
\\
&=
\|H\|_{\infty,2}.
\end{aligned}
\]

On obtient donc

\[
\boxed{
\|\widetilde A H\|_{\infty,2}
\leq
\|H\|_{\infty,2}.
}
\]

La même démonstration appliquée à \(H-K\) donne

\[
\boxed{
\|\widetilde A H-\widetilde A K\|_{\infty,2}
\leq
\|H-K\|_{\infty,2}.
}
\]

L'opérateur relationnel est donc non expansif pour cette norme.

\section{Bornes et régularité de la saturation}
\label{annexe-existence-tanh}

La fonction \(\tanh\) est appliquée composante par composante.

Pour tout \(y\in\mathbb R^d\),

\[
|\tanh(y_k)|\leq1,
\]

donc

\[
\boxed{
\|\tanh(y)\|_2\leq\sqrt d.
}
\]

Par ailleurs,

\[
|\tanh'(s)|=\operatorname{sech}^2(s)\leq1.
\]

Le théorème des accroissements finis, appliqué à chaque composante,
donne

\[
|\tanh(y_k)-\tanh(z_k)|
\leq
|y_k-z_k|.
\]

Après sommation,

\[
\boxed{
\|\tanh(y)-\tanh(z)\|_2
\leq
\|y-z\|_2.
}
\]

En combinant ce résultat avec la non-expansivité de
\(\widetilde A\),

\[
\boxed{
\left\|
\tanh(\widetilde A H)
-
\tanh(\widetilde A K)
\right\|_{\infty,2}
\leq
\|H-K\|_{\infty,2}.
}
\]

\section{Solution et borne de la trace}
\label{annexe-existence-trace}

Pour chaque nœud,

\[
\dot\tau_i
=
\alpha\,\tanh\!\bigl((\widetilde A H)_i\bigr)
-
\beta\tau_i,
\qquad
\beta>0.
\]

Entre deux instants \(t_0\leq t\), la formule de variation des
constantes donne

\[
\tau_i(t)
=
e^{-\beta(t-t_0)}\tau_i(t_0)
+
\alpha
\int_{t_0}^t
e^{-\beta(t-s)}
\tanh\!\bigl((\widetilde A H(s))_i\bigr)
\,\mathrm ds.
\]

En utilisant la borne précédente,

\[
\begin{aligned}
\|\tau_i(t)\|_2
&\leq
e^{-\beta(t-t_0)}
\|\tau_i(t_0)\|_2
\\
&\quad+
\alpha\sqrt d
\int_{t_0}^t
e^{-\beta(t-s)}
\,\mathrm ds
\\
&=
e^{-\beta(t-t_0)}
\|\tau_i(t_0)\|_2
+
\frac{\alpha\sqrt d}{\beta}
\left(
1-e^{-\beta(t-t_0)}
\right).
\end{aligned}
\]

Le membre de droite est une combinaison convexe de
\(\|\tau_i(t_0)\|_2\) et de \(\alpha\sqrt d/\beta\). Ainsi,

\[
\boxed{
\|\tau_i(t)\|_2
\leq
\max\left(
\|\tau_i(t_0)\|_2,
\frac{\alpha\sqrt d}{\beta}
\right).
}
\]

Comme la trace ne change pas lors d'un Kairos, cette borne reste
valable le long de toute l'exécution hybride.

\section{Préconditionneur : formule et bornes spectrales}
\label{annexe-existence-preconditionneur}

On pose

\[
Q(\tau)
=
I_d+
\frac{\lambda}{1+\|\tau\|^2}
\tau\tau^\top,
\qquad
\lambda\geq0.
\]

Si \(z\perp\tau\), alors

\[
Q(\tau)z=z.
\]

Dans la direction de \(\tau\),

\[
Q(\tau)\tau
=
\left(
1+
\frac{\lambda\|\tau\|^2}
{1+\|\tau\|^2}
\right)\tau.
\]

Les valeurs propres de \(Q(\tau)\) sont donc

\[
1
\]

dans les directions orthogonales à \(\tau\), et

\[
q_\parallel(\tau)
=
1+
\frac{\lambda\|\tau\|^2}
{1+\|\tau\|^2}
\]

dans la direction de \(\tau\). En particulier,

\[
1\leq q_\parallel(\tau)<1+\lambda.
\]

La matrice \(Q(\tau)\) est symétrique définie positive et donc
inversible.

En appliquant la formule de Sherman--Morrison
\cite{sherman1950adjustment}, on obtient

\[
\boxed{
M(\tau)
=
Q(\tau)^{-1}
=
I_d
-
\frac{\lambda}
{1+(1+\lambda)\|\tau\|^2}
\tau\tau^\top.
}
\]

Les valeurs propres de \(M(\tau)\) sont \(1\) dans les directions
orthogonales à \(\tau\), et

\[
m_\parallel(\tau)
=
\frac{1+\|\tau\|^2}
{1+(1+\lambda)\|\tau\|^2}
\]

dans la direction de \(\tau\).

Comme

\[
\frac{1}{1+\lambda}
\leq
m_\parallel(\tau)
\leq1,
\]

on a, pour tout \(y\in\mathbb R^d\),

\[
\boxed{
m_0\|y\|_2^2
\leq
y^\top M(\tau)y
\leq
\|y\|_2^2,
\qquad
m_0:=\frac{1}{1+\lambda}.
}
\]

En notation matricielle,

\[
\boxed{
m_0I_d
\preceq
M(\tau)
\preceq
I_d.
}
\]

Il en résulte notamment

\[
\boxed{
\|M(\tau)\|_{\mathrm{op}}\leq1.
}
\]

\section{Régularité et dérivée de \texorpdfstring{$M$}{M}}
\label{annexe-existence-derivee-M}

Posons

\[
c=1+\lambda,
\qquad
a(\tau)
=
\frac{\lambda}{1+c\|\tau\|^2}.
\]

Alors

\[
M(\tau)=I_d-a(\tau)\tau\tau^\top.
\]

Pour une direction \(h\in\mathbb R^d\),

\[
Da(\tau)[h]
=
-\frac{2\lambda c\,\tau\cdot h}
{\left(1+c\|\tau\|^2\right)^2}.
\]

Par conséquent,

\[
\begin{aligned}
DM(\tau)[h]
={}&
\frac{2\lambda c\,(\tau\cdot h)}
{\left(1+c\|\tau\|^2\right)^2}
\tau\tau^\top
\\
&-
\frac{\lambda}
{1+c\|\tau\|^2}
\left(
h\tau^\top+\tau h^\top
\right).
\end{aligned}
\]

Cette expression est définie et continue pour tout \(\tau\), y
compris en \(\tau=0\). La fonction \(M\) est donc lisse.

Nous pouvons également obtenir une borne globale sur sa différentielle.
En posant \(r=\|\tau\|_2\),

\[
\left\|
\frac{\lambda}
{1+cr^2}
\left(
h\tau^\top+\tau h^\top
\right)
\right\|_{\mathrm{op}}
\leq
\frac{2\lambda r}{1+cr^2}\|h\|_2.
\]

Or

\[
\sup_{r\geq0}\frac{r}{1+cr^2}
=
\frac{1}{2\sqrt c}.
\]

Le premier terme vérifie

\[
\left\|
\frac{2\lambda c\,(\tau\cdot h)}
{(1+cr^2)^2}
\tau\tau^\top
\right\|_{\mathrm{op}}
\leq
\frac{2\lambda c r^3}{(1+cr^2)^2}
\|h\|_2.
\]

Après le changement de variable \(y=\sqrt c\,r\), on a
\(r^3=c^{-3/2}y^3\) et \((1+cr^2)^2=(1+y^2)^2\). Le facteur
\(c\,c^{-3/2}=c^{-1/2}\) donne donc

\[
\sup_{r\geq0}
\frac{2\lambda c r^3}{(1+cr^2)^2}
=
\frac{2\lambda}{\sqrt c}
\sup_{y\geq0}
\frac{y^3}{(1+y^2)^2}.
\]

La fonction \(y^3/(1+y^2)^2\) atteint son maximum en
\(y=\sqrt3\), avec

\[
\frac{y^3}{(1+y^2)^2}
=
\frac{3\sqrt3}{16}.
\]

On obtient finalement

\[
\boxed{
\|DM(\tau)[h]\|_{\mathrm{op}}
\leq
L_M(\lambda)\|h\|_2,
}
\]

avec la constante explicite

\[
\boxed{
L_M(\lambda)
=
\frac{\lambda}{\sqrt{1+\lambda}}
\left(
1+\frac{3\sqrt3}{8}
\right).
}
\]

Cette borne est conservative, mais elle montre que \(M\) est
globalement lipschitzienne.

Dans la suite, on utilise la borne uniforme
\[
\bar\tau
:=
\max_i\max\left\{\|\tau_i(t_0)\|_2,
\frac{\alpha\sqrt d}{\beta}\right\},
\]
qui est précisément celle obtenue à la section précédente.

Comme

\[
\|\dot\tau_i\|_2
\leq
\alpha\sqrt d+\beta\bar\tau,
\]

elle donne aussi

\[
\boxed{
\|\dot M(\tau_i)\|_{\mathrm{op}}
\leq
L_M(\lambda)
\left(
\alpha\sqrt d+\beta\bar\tau
\right).
}
\]

C'est la constante notée \(\mu\) dans l'annexe consacrée à l'énergie.

\section{Régularité de la friction}
\label{annexe-existence-friction}

La friction est

\[
\Gamma(\tau)
=
\gamma_0+\gamma_1\|\tau\|_2.
\]

L'inégalité triangulaire inverse donne

\[
\left|
\|\tau\|_2-\|\sigma\|_2
\right|
\leq
\|\tau-\sigma\|_2.
\]

Ainsi,

\[
\boxed{
|\Gamma(\tau)-\Gamma(\sigma)|
\leq
\gamma_1\|\tau-\sigma\|_2.
}
\]

La fonction \(\Gamma\) est donc globalement lipschitzienne. Elle n'est
pas différentiable en \(\tau=0\) lorsque \(\gamma_1>0\), mais la
différentiabilité n'est pas nécessaire à l'existence et à l'unicité
des solutions.

\section{Champ de vecteurs pour un graphe fixé}
\label{annexe-existence-champ}

Pour un graphe \(A\) fixé, la dynamique continue s'écrit

\[
\dot z=F_A(t,z).
\]

La seule dépendance explicite en \(t\) vient de l'entrée mesurable
\(u(t)\).

Nous vérifions les hypothèses de Carathéodory.

\subsection{Mesurabilité en temps}
\label{annexe-existence-mesurabilite}

Pour tout état \(z\) fixé, les termes internes sont constants en
\(t\) : \(M(\tau)\), \(\Gamma(\tau)\) et
\(\tanh(\widetilde A H)\) sont des fonctions de l'état seul. La
dépendance temporelle explicite est entièrement portée par l'entrée
mesurable \(u(t)\). La fonction

\[
t\longmapsto F_A(t,z)
\]

est donc mesurable.

\subsection{Continuité et caractère localement lipschitzien}
\label{annexe-existence-lipschitz}

Pour presque tout \(t\), les termes suivants sont localement
lipschitziens en \(z\) :

\begin{itemize}
\item
  \(V\mapsto V\) ;
\item
  \(H\mapsto\tanh(\widetilde A H)\), qui est même globalement
  lipschitzien ;
\item
  \((\tau,V)\mapsto\Gamma(\tau)V\) ;
\item
  \((\tau,H)\mapsto M(\tau)H\) ;
\item
  \((H,V,P)\mapsto
  \kappa_1\|H_i\|+\kappa_2\|V_i\|-\kappa_3P_i\).
\end{itemize}

À titre d'exemple, sur un compact où
\(\|V\|_{\infty,2}\leq R_V\) et
\(\|\tau\|_{\infty,2}\leq R_\tau\),

\[
\begin{aligned}
\|\Gamma(\tau_i)V_i-\Gamma(\sigma_i)W_i\|
\leq{}&
\Gamma_{\max}\|V_i-W_i\|
\\
&+
\gamma_1R_V\|\tau_i-\sigma_i\|,
\end{aligned}
\]

où

\[
\Gamma_{\max}
=
\gamma_0+\gamma_1R_\tau.
\]

De même,

\[
\begin{aligned}
\|M(\tau_i)H_i-M(\sigma_i)K_i\|
\leq{}&
\|H_i-K_i\|
\\
&+
L_M(\lambda)R_H\|\tau_i-\sigma_i\|
\end{aligned}
\]

sur un compact où \(\|K_i\|\leq R_H\).

Le champ \(F_A(t,\cdot)\) est donc localement lipschitzien.

\subsection{Majoration sur les compacts}
\label{annexe-existence-majoration-compacte}

Sur tout compact de l'espace des états, les variables \(H\), \(V\),
\(\tau\) et \(P\) sont bornées. Les fonctions \(\tanh\), \(M\) et
\(\Gamma\) y sont bornées. L'entrée \(u(t)\) est essentiellement
bornée.

Il existe donc une fonction intégrable \(m_K(t)\), que l'on peut ici
choisir constante, telle que

\[
\|F_A(t,z)\|
\leq
m_K(t)
\]

pour presque tout \(t\) et tout \(z\) dans le compact \(K\).

\section{Existence et unicité locale}
\label{annexe-existence-locale}

Les trois propriétés précédentes sont les hypothèses usuelles du
théorème de Carathéodory
\cite{filippov1988differential}.

\begin{quote}
\resultheading{Proposition}{prop:A-existence-unicite-locale}{}
Pour tout graphe fixé \(A\), toute condition initiale continue
\(z(t_0)=z_0\) et toute entrée mesurable essentiellement bornée, il
existe une solution locale absolument continue. Cette solution est
unique et satisfait les équations presque partout.
\end{quote}

La solution peut être prolongée tant qu'elle reste dans un ensemble
borné. Il reste donc à exclure une explosion en temps fini.

\section{Absence d'explosion en temps fini}
\label{annexe-existence-non-explosion}

Posons

\[
B=\zeta\sqrt d+u_{\max}.
\]

Pour une fonction absolument continue, la norme est absolument
continue. Aux instants où \(H_i\neq0\),

\[
\frac{\mathrm d}{\mathrm dt}\|H_i\|_2
=
\frac{H_i}{\|H_i\|_2}\cdot V_i
\leq
\|V_i\|_2.
\]

Plus généralement, pour toute fonction absolument continue, la
sous-additivité de la norme donne partout l'inégalité en dérivée
supérieure de Dini

\[
D^+\|H_i\|_2\leq\|V_i\|_2,
\]
ce qui inclut sans exception le cas \(H_i=0\).

Lorsque \(V_i\neq0\),

\[
\begin{aligned}
\frac{\mathrm d}{\mathrm dt}\|V_i\|_2
&=
\frac{V_i}{\|V_i\|_2}\cdot\dot V_i
\\
&=
-\Gamma(\tau_i)\|V_i\|_2
-
\frac{V_i}{\|V_i\|_2}\cdot M(\tau_i)H_i
\\
&\quad+
\zeta
\frac{V_i}{\|V_i\|_2}
\cdot
\tanh\!\bigl((\widetilde A H)_i\bigr)
+
\frac{V_i}{\|V_i\|_2}\cdot u_i(t).
\end{aligned}
\]

Comme \(\Gamma\geq0\), \(\|M(\tau_i)\|_{\mathrm{op}}\leq1\) et
\(\|\tanh(\cdot)\|_2\leq\sqrt d\),

\[
D^+\|V_i\|_2
\leq
\|H_i\|_2+B.
\]

Définissons

\[
Y(t)
=
\max_i
\left(
\|H_i(t)\|_2+\|V_i(t)\|_2
\right).
\]

La dérivée supérieure du maximum vérifie

\[
D^+Y(t)
\leq
Y(t)+B.
\]

L'inégalité de Grönwall donne, pour \(t\geq t_0\),

\[
\boxed{
Y(t)
\leq
\left(
Y(t_0)+B
\right)
e^{t-t_0}
-
B.
}
\]

Cette estimation est finie sur tout intervalle de temps fini.

La trace est bornée par la formule obtenue plus haut. La pression
vérifie

\[
\dot P_i
+
\kappa_3P_i
=
\kappa_1\|H_i\|_2+\kappa_2\|V_i\|_2.
\]

Ainsi,

\[
\begin{aligned}
P_i(t)
={}&
e^{-\kappa_3(t-t_0)}P_i(t_0)
\\
&+
\int_{t_0}^t
e^{-\kappa_3(t-s)}
\left(
\kappa_1\|H_i(s)\|_2
+
\kappa_2\|V_i(s)\|_2
\right)
\,\mathrm ds.
\end{aligned}
\]

Le membre de droite est fini sur tout intervalle borné. Il est
également non négatif si \(P_i(t_0)\geq0\).

Aucune composante continue ne peut donc diverger en temps fini.

\begin{quote}
\resultheading{Proposition}{prop:A-non-explosion-flot}{}
Pour un graphe fixé, la solution de Carathéodory unique se prolonge
pour tout temps positif, tant qu'aucun saut n'est imposé par la règle
hybride.
\end{quote}

Cette affirmation ne signifie pas que les trajectoires restent
uniformément bornées lorsque \(t\to\infty\). Elle exclut seulement une
explosion à temps fini.

\section{Premier instant de saut}
\label{annexe-existence-premier-saut}

Supposons l'état initial admissible :

\[
0\leq P_i(t_0)\leq\Theta
\]

pour tout \(i\).

Chaque \(P_i(t)\) est continu pendant le flot. Deux cas sont possibles.

\begin{enumerate}
\item
  Aucun \(P_i\) n'atteint \(\Theta\). La solution continue se poursuit
  indéfiniment.
\item
  Au moins une pression atteint \(\Theta\). On définit
  \[
  t_1
  =
  \inf\left\{
  t\geq t_0:
  \max_iP_i(t)=\Theta
  \right\}.
  \]
  Par continuité, ce temps est atteint et au moins un indice vérifie
  \(P_i(t_1)=\Theta\).
\end{enumerate}

La règle de priorité impose alors un saut.

\section{Le saut est bien défini}
\label{annexe-existence-saut-bien-defini}

À l'instant du saut, le modèle choisit

\[
i^\ast
=
\min\left\{
i:
P_i=\Theta
\right\}.
\]

Les similarités sont définies par

\[
s(i,j)
=
\frac{\langle\tau_i,\tau_j\rangle}
{\|\tau_i\|_2\|\tau_j\|_2}
\]

lorsque les deux traces sont non nulles, et par \(s(i,j)=0\) sinon.

Les égalités de similarité sont départagées par l'indice. Les ensembles
de voisins et de non-voisins sont finis. Le choix des arêtes à retirer
et à ajouter est donc déterministe.

Le nombre d'arêtes modifiées est

\[
m
=
\min\left(
2,
|\mathcal N(i^\ast)|,
|\mathcal N^c(i^\ast)|
\right).
\]

Le saut retire \(m\) arêtes existantes et ajoute \(m\) arêtes
absentes. Il ne crée ni boucle ni arête multiple. La symétrie de la
matrice d'adjacence est maintenue.

Ainsi :

\[
\boxed{
\text{le graphe reste simple, non orienté et conserve son nombre
d'arêtes.}
}
\]

L'ajout de la boucle de calcul \(I\) dans la définition de
\(\widetilde A\) garantit que la normalisation reste définie après le
recâblage, même si un nœud n'a plus aucun voisin.

Les autres composantes vérifient

\[
H^+=H^-,
\qquad
\tau^+=\tau^-,
\]

\[
V_{i^\ast}^+
=
r_VV_{i^\ast}^-,
\qquad
P_{i^\ast}^+
=
r_P\Theta,
\]

avec \(r_V,r_P\in[0,1)\).

Le saut ne peut donc pas augmenter
\(\|H\|_{\infty,2}\), \(\|\tau\|_{\infty,2}\) ni
\(\|V\|_{\infty,2}\). Il préserve également \(P_i\geq0\).

\section{Événements simultanés}
\label{annexe-existence-evenements-simultanes}

Plusieurs pressions peuvent atteindre \(\Theta\) au même instant. Les
sauts sont alors appliqués successivement sans faire avancer le temps
continu.

Après le saut du nœud sélectionné,

\[
P_{i^\ast}^+
=
r_P\Theta
<
\Theta.
\]

Les pressions des autres nœuds ne changent pas. Le nœud sélectionné
sort donc de l'ensemble des indices au seuil, tandis qu'aucun nouvel
indice n'y entre sans écoulement du temps.

Il y a par conséquent au plus \(N\) sauts dans une cascade
instantanée.

\section{Borne sur un horizon fini malgré les sauts}
\label{annexe-existence-borne-horizon-fini}

La borne de Grönwall établie pour le flot reste valable le long d'une
exécution hybride.

En effet, lors d'un saut,

\[
H^+=H^-,
\]

et

\[
\|V_{i^\ast}^+\|_2
=
r_V\|V_{i^\ast}^-\|_2
\leq
\|V_{i^\ast}^-\|_2.
\]

La quantité

\[
Y
=
\max_i
\left(
\|H_i\|_2+\|V_i\|_2
\right)
\]

ne peut donc pas augmenter lors d'un saut. Les changements du graphe
ne modifient pas la constante \(B\), car la matrice normalisée reste
stochastique par lignes après chaque recâblage.

Pour tout horizon continu \(T>0\),

\[
\boxed{
Y(t,j)
\leq
\left(
Y(0,0)+B
\right)e^T-B
=:
R_T
}
\]

pour tout point hybride \((t,j)\) tel que \(0\leq t\leq T\).

Cette borne dépend de \(T\), mais elle ne dépend pas du nombre de
sauts déjà effectués.

\section{Absence d'accumulation de sauts en temps fini}
\label{annexe-existence-non-zenon}

Fixons un horizon continu \(T>0\). La borne précédente donne, pour
tout nœud et tout \(t\leq T\),

\[
\|H_i(t)\|_2\leq R_T,
\qquad
\|V_i(t)\|_2\leq R_T.
\]

Ainsi,

\[
\dot P_i
\leq
\kappa_1R_T+\kappa_2R_T
=:
K_T.
\]

Si \(K_T=0\), aucune pression ne peut remonter après un saut.

Supposons \(K_T>0\). Après un saut du nœud \(i\),

\[
P_i^+=r_P\Theta.
\]

Pour atteindre à nouveau \(\Theta\), il doit gagner au moins

\[
(1-r_P)\Theta.
\]

Le temps continu séparant deux sauts successifs du même nœud est donc
au moins

\[
\boxed{
\delta_T
=
\frac{(1-r_P)\Theta}{K_T}
>0.
}
\]

Sur l'intervalle \([0,T]\), chaque nœud saute au plus

\[
1+\frac{T}{\delta_T}
\]

fois, en comptant un éventuel saut initial. Le nombre total de sauts
est donc majoré par

\[
\boxed{
N\left(
1+\frac{T}{\delta_T}
\right).
}
\]

Il est fini.

\begin{quote}
\resultheading{Proposition}{prop:A-absence-zenon}{}
Sous les hypothèses du noyau, avec \(r_P<1\), une exécution issue d'un
état admissible ne peut pas contenir une infinité de sauts sur un
intervalle de temps continu fini.
\end{quote}

Cette propriété exclut le phénomène de Zénon
\cite{goebel2012hybrid}.

\section{Existence et unicité de l'exécution hybride}
\label{annexe-existence-hybride}

Nous pouvons maintenant assembler les résultats.

\begin{enumerate}
\item
  Pour un graphe donné, le flot possède une solution unique.
\item
  Le premier instant de saut est déterminé par la continuité des
  pressions.
\item
  Le choix du nœud et le recâblage sont déterministes grâce aux règles
  de départage.
\item
  Une cascade simultanée comporte au plus \(N\) sauts.
\item
  Il n'y a pas d'accumulation infinie de sauts en temps fini.
\item
  Entre deux sauts, aucune composante ne peut exploser en temps fini.
\end{enumerate}

\begin{quote}
\resultheading{Proposition}{prop:A-execution-hybride-unique}{}
Pour toute condition initiale admissible et toute entrée mesurable
essentiellement bornée, le noyau définit une exécution hybride unique
pour tout temps continu positif.
\end{quote}

L'unicité est entendue pour les conventions déterministes retenues :
priorité au saut sur \(C\cap D\), choix du plus petit indice et
départage des similarités par l'indice.

\section{Portée et limites}
\label{annexe-existence-limites}

Les résultats de cette annexe établissent :

\[
\boxed{
\text{existence, unicité et absence de Zénon sur tout horizon fini.}
}
\]

Ils ne montrent pas que les trajectoires \(H(t)\) et \(V(t)\) restent dans un même ensemble borné lorsque \(t\to\infty\).

Cette question asymptotique est différente. Elle est traitée par les estimations d'énergie de l'Annexe~\ref{annexe-energie-flots-sauts}, sous des conditions suffisantes explicites.

Enfin, le saut est déterministe mais il n'est pas continu par rapport à l'état aux points où l'ordre des similarités change. La robustesse des trajectoires vis-à-vis de petites perturbations des classements n'est pas étudiée ici.
